\begin{prop}


Отображение $\varphi:A\to A^{\prime}$, определенное на измеримом


множестве $A\subset\Bbb R^n$ и имеющее


аппроксимативные частные производные на множестве $A$, порождает


ограниченный оператор вложения


$$


\varphi^{\ast} :


L_r(A^{\prime})\to L_s(A), \quad 1\leq s\leq r<\infty,


$$


тогда и только тогда, когда $\varphi$ обладает ${\cal N}^{-1}$-свойством,


его якобиан $J(x,\varphi)$ почти всюду отличен от нуля


и объемная производная


$$


\big(J_{\varphi^{-1}}(y)\big)^{\frac{1}{s}}=


\bigg(\,\sum\limits_{x\in\varphi^{-1}(y)\setminus\Sigma_{\varphi}}


\frac{1}{|J(x,\varphi)|}\bigg)^{\frac{1}{s}}\in L_{\lambda}(A^{\prime}),


$$


где $\lambda=\frac{rs}{r-s}$ при $s<r$ и $\lambda=\infty$ при $s=r$.


Норма оператора $\varphi^{\ast}$ равна


$\|(J_{\varphi^{-1}})^{\frac{1}{s}}\mid


L_\lambda(A^{\prime})\|$.


\end{prop}





\begin{pf}


Докажем только невырожденность якобиана.


${\cal N}^{-1}$-свойство отображения $\varphi$ эквивалентно


тому, что образ множества $B\subset A$, $B\cap \Sigma_{\varphi}=\emptyset$,


положительной меры есть множество положительной меры. Пусть


$Z=\{x\in A\setminus \Sigma_{\varphi}\mid J(x,\varphi)=0\}$.


Если $|Z|>0$, то $|\varphi(Z\setminus \Sigma_{\varphi})|>0$.


Однако по формуле (1) $|\varphi(Z\setminus \Sigma_{\varphi})|=0$.


Полученное противоречие приводит к тому, что $|Z|=0$.


\end{pf}





Обозначим символом $W^1_{s,q}(D)$, $s\in[1,\infty]$, пространство


функций из $L^1_{q}(D)$, имеющих конечную норму


$\|f\mid W^1_{s,q}(D)\|=\|f\mid L_s(D)\|+\|\nabla f\mid L_q(D)\|$.





Из теоремы 1 и предложения 4 получается





\begin{thm}


Пусть непостоянное отображение


$\varphi : D\to D^{\prime}$ порождает ограниченный


оператор вложения


$$


\varphi^{\ast} : L^1_p(D^{\prime})\to


L^1_q(D), \quad 1\leq q\leq p<n.


$$


Тогда


\begin{itemize}


\item[1)] объемная производная $J_{\varphi^{-1}}$ {\em ``}обратного


отображения{\em ''}


принадлежит $L_t(D^{\prime})$, где


$t=\frac{p}{p-q}\frac{n-q}{n};$


\item[2)] на всяком измеримом множестве $A^{\prime}\subset D^{\prime}$


отображение $\varphi$ порождает по правилу суперпозиции


$\varphi^{\ast} f = f\circ\varphi$ ограниченный оператор вложения


$$


\varphi^{\ast}: L_r(A^{\prime})\to L_s(\varphi^{-1}(A^{\prime})),


\ \ \text{где} \ \ \frac{p(n-q)}{q(n-p)}\leq r<\infty


\ \ \text{и}\ \ s=r \frac{q(n-p)}{p(n-q)};


$$


\item[3)] норма оператора


$\varphi^{\ast}$


удовлетворяет соотношению


$$


\|\varphi^{\ast}:L_r(A^{\prime})\to L_s(\varphi^{-1}(A^{\prime}))\|\leq


\|H_q(\cdot)\mid L_\varkappa(A^{\prime})\|^{\frac{pn}{r(n-p)}};


$$


\item[4)] отображение $\varphi$ обладает ${\cal N}^{-1}$-свойством,


и его якобиан $J(x,\varphi)$ почти всюду отличен от нуля\/$;$


\item[5)]


отображение $\varphi$ непостоянно на каждой компоненте связности


открытого множества $D;$


\item[6)]


если $r\geq \frac{p(n-q)}{q(n-p)}$, то отображение


$\varphi$ порождает ограниченный оператор вложения


$$


\varphi^{\ast}: W^1_{r,p}(D^{\prime})\to W^1_{s,q}(D),


\ \ \text{где} \ \ s=r \frac{q(n-p)}{p(n-q)}.


$$


\end{itemize}


В частности, если


множество $D$ ограничено при $q<p<n$, и оператор вложения


$i:W^1_p(D^{\prime})\to L_{p^*}(D^{\prime})$ ограничен,


$p^*=\frac{pn}{n-p}$, то отображение $\varphi$ порождает


ограниченный оператор вложения $\varphi^{\ast}:


W^1_p(D^{\prime})\to W^1_{q}(D)$.


\end{thm}





\begin{pf}


Согласно теоремe 1 функция $H_q(y)$,


определенная формулой (2), принадлежит пространству


$L_\varkappa(D^{\prime})$. Для почти всех $y\in A^{\prime}\subset


D^{\prime}$ имеем


\begin{multline*}


H_q(y)= \bigg(\,\sum_{x\in\varphi^{-1}(y)\setminus


\Sigma_{\varphi}}


\frac{|D\varphi|^q(x)}{|J(x,\varphi)|}\bigg)^{\frac{1}{q}}=


\bigg(\,\sum_{x\in\varphi^{-1}(y)\setminus \Sigma_{\varphi}}


\bigg(\frac{|D\varphi|^n(x)}{J(x,\varphi)|}\bigg)^{\frac{q}{n}}


\frac{1}{|J(x,\varphi)|^{1-\frac{q}{n}}}\bigg)^{\frac{1}{q}}\ge \\


%


\geq \bigg(\,\sum_{x\in\varphi^{-1}(y)\setminus


\Sigma_{\varphi}} \frac{1}{|J(x,\varphi)|^{1-\frac{q}{n}}}


\bigg)^{\frac{1}{q}}\geq


\bigg(\,\sum_{x\in(\varphi^{-1}(y)\cap A)\setminus


\Sigma_{\varphi}} \frac{1}{|J(x,\varphi)|}


\bigg)^{\frac{1}{q}-\frac{1}{n}}


=J_{\varphi^{-1}}(y)^{\frac{1}{q}-\frac{1}{n}}


\end{multline*}


в силу поточечного соотношения


$|J(x,\varphi)|\leq |D\varphi|^n(x)$ и неравенства Йенсена.


Отсюда получим


\begin{alignat*}{2}


\|J_{\varphi^{-1}}(\cdot)\mid L_t(A^{\prime})\|^t&\leq


\|H_q(\cdot)\mid L_\varkappa(A^{\prime})\|^\varkappa&\quad &\text{при}


\ \ q<p \\


\intertext{и}


\esssup_{y\in


A^{\prime}}J_{\varphi^{-1}}(y)^{\frac{1}{q}-\frac{1}{n}}&\leq


\esssup_{y\in A^{\prime}}H_q(y)&\quad &\text{при}


\ \ q=p.


\end{alignat*}


Из этих оценок объемной производной ``обратного


отображения'' и предложения 4 получим, что отображение $\varphi$


порождает ограниченный оператор вложения


$$


\varphi^{\ast}:


L_r(A^{\prime})\to L_s(\varphi^{-1}(A^{\prime})),\quad \text{где}


\ \ \frac{p(n-q)}{q(n-p)}\leq r<\infty,\ \ \text{а} \ \ s=r


\frac{q(n-p)}{p(n-q)}.


$$


Из предложения 4 получим ${\cal N}^{-1}$-свойство отображения


$\varphi$ и невырожденность якобиана почти всюду.





%% Из ${\cal


%% N}^{-1}$-свойства отображения $\varphi$ вытекает его непостоянство


%% на каждой компоненте связности открытого множества $D$.





Далее из соотношений


$$


\|(J_{\varphi^{-1}}(\cdot))^{\frac{1}{s}}\mid


L_\lambda(A^{\prime})\|^\lambda=


\|J_{\varphi^{-1}}(\cdot)\mid


L_t(A^{\prime})\bigr\|^t\leq


\|H_q(\cdot)\mid L_\varkappa(A^{\prime})\|^\varkappa,\quad q<p,


$$


получаем оценку для нормы оператора


$\varphi^{\ast}: L_r(D^{\prime})\to L_s(D)$,


$$


\|\varphi^{\ast}:L_r(A^{\prime})\to L_s(\varphi^{-1}(A^{\prime}))\|\leq


\|H_q(\cdot)\mid L_\varkappa(A^{\prime})\|^{\frac{\varkappa}{\lambda}}=


\|H_q(\cdot)\mid L_\varkappa(A^{\prime})\|^{\frac{pn}{r(n-p)}},


$$


справедливую и при $q=p$.





Учитывая вышесказанное, получим ограниченность оператора суперпозиции в


пространствах Соболева с основной нормой.


\end{pf}





Из теоремы 3 вытекает утверждение, обобщающее соответствующий


результат из [18].





\begin{cor}


Oтображение $\varphi : D\to


D^{\prime}$ порождает ограниченный оператор вложения


$\varphi^{\ast}: W^1_p(D^{\prime})\to W^1_{p}(D)$,


$1\leq p<n$, тогда и только тогда, когда оно порождает


ограниченный оператор вложения


$\varphi^{\ast} : L^1_p(D^{\prime})\to L^1_p(D)$.


\end{cor}





\begin{pf}


Достаточность доказана в предыдущем утверждении. Чтобы установить


необходимость, надо доказать, что выполняются условия теоремы 1, если


оператор $\varphi^{\ast}: W^1_p(D^{\prime})\to W^1_{p}(D)$ ограничен.


Свойство $\ACL$ отображения $\varphi$ доказано в предложении 2.





Фиксируем срезку $\eta\in C_{0}^{\infty}({\Bbb R}^n)$, равную


единице на $B(0,1)$ и нулю вне шара $B(0,2)$. Подставим в


неравенство $\|\nabla\varphi^{\ast}(f)\mid L_p(D)\|


\leq\|\varphi^{\ast}\|\|f\mid W_p^1(D^{\prime})\|$ вместо $f$


функции


$h_{j}(z)=(z-y)_{j}\eta(\frac{z-y}{r})$, $j=1,\dotsc,n$, где


$(z-y)_{j}$ обозначает $j$-ю координату вектора $z-y$,


$B(y,2r)\subset D^{\prime}$. Получим


$$


\int\limits_{\varphi^{-1}(B(y,r))}|D\varphi|^p(x)\,dx \leq


C\|\varphi^{\ast}\|^p|B(y,2r)|,\quad r\in(0,1),


$$


поскольку


$\|h_{j} \mid L_p(D^{\prime})\|^p=o(|B(y,2r)|)$ равномерно по


всем шарам $B(y,2r)\subset D^{\prime}$ и $r\in(0,1)$ (здесь


$C$ --- некоторая постоянная, зависящая только от размерности $n$


и показателя $p$). Далее так же, как и в доказательстве теоремы 1,


начиная с формулы (3), показывается, что отображение


$\varphi$ имеет конечное искажение, и справедлива оценка


$$


\bigg(\,\sum\begin{Sb} x\in\varphi^{-1}(y)\setminus


\Sigma_{\varphi},\\ J(x,\varphi)\ne0\end{Sb}


\frac{|D\varphi|^p(x)}{|J(x,\varphi)|}\bigg)^{\frac{1}{p}} \leq


C\|\varphi^*\|


$$


для почти всех $y\in D^{\prime}$ таких, что


$\varphi^{-1}(y)\setminus (\Sigma_{\varphi}\cup Z)\ne\emptyset$,


и $J(x,\varphi)\ne0$,


если $x\in \varphi^{-1}(y)\setminus (\Sigma_{\varphi}\cup Z)$.


По теореме 1 оператор $\varphi^{\ast} : L^1_p(D^{\prime})\to L^1_p(D)$


ограничен.


\end{pf}








\section{Оценки искажения меры}





Рассмотрим более детально вопросы, связанные со свойствами искажения меры.





\begin{thm}


Если отображение $\varphi:D\to D^{\prime}$


порождает ограниченный оператор вложения


$$


\varphi^{\ast} :


L^1_p(D^{\prime})\to L^1_q(D),\quad 1\leq q\leq p\leq n,


$$


то $\varphi$ обладает ${\cal N}^{-1}$-свойством.


\end{thm}





\begin{pf}


Утверждение теоремы при $p<n$ доказано в теореме 3.


Здесь мы приведем другое доказательство, единое для всех показателей


$p\leq n$.


%% Поскольку точка имеет нулевую $n$-емкость, то отображение


%% $\varphi$ не может быть постоянным ни


%% на какой компоненте связности открытого множества $D$.





По теореме 1 отображение $\varphi\in\ACL(D)$, и можно


считать, что открытое множество $D$ ограничено. Далее


рассмотрим случай $q<n$; случай $q=n$ сводится к


предыдущему, поскольку множество $D$ ограниченное.





Фиксируем срезку $\eta\in C_{0}^{\infty}({\Bbb R}^n)$, равную


единице на $B(0,1)$ и нулю вне шара $B(0,2)$. На основании леммы~1


имеем


$$


\|\varphi^{\ast}f\mid L^1_q(D)\| \leq


C_1{\Phi(2B)}^{\frac{1}{\sigma}} |B|^{\frac{1}{p} - \frac{1}{n}},


$$


где $B=B(y_0,r)$ --- такой шар, что $\ov{B(y_0,2r)}\subset D'$,


и $f(y)=\eta\big(\frac{y-y_0}{r}\big)$. (В случае $p=q$ полагаем


$\Phi(2B)^{\frac{1}{\sigma}}=\|\varphi^{\ast}\|$ для всех шаров $B$.)


Возьмем в $D'$ произвольное множество $E$ нулевой меры.


Без ограничения общности можно считать, что множество $E$


ограничено. Так как отображение $\varphi$ имеет конечное


искажение, то $\varphi^{-1}(E)\ne D$ (в противном случае


$J(x,\varphi)=0$ и, следовательно, $D\varphi(x)=0$, откуда


получаем, что $\varphi$ --- постоянное отображение).


Поэтому найдется куб $Q\subset D$ такой, что $2Q\subset D$ и


$|Q\setminus\varphi^{-1}(E)|>0$ (здесь $2Q$ --- куб с тем же


центром, что и $Q$, с ребрами, растянутыми в два раза сравнительно


с $Q$). Поскольку компоненты данного отображения измеримы, то по


теореме Лузина найдется компакт $T\subset


Q\setminus\varphi^{-1}(E)$ положительной меры такой, что $\varphi


: T\to D^{\prime}$ непрерывно. Тогда образ $\varphi(T)\subset


D^{\prime}$ компактен и $\varphi(T)\cap E=\emptyset$. Рассмотрим


произвольное открытое множество $U\supset E$, $\varphi(T)\cap


U=\emptyset$, $U\subset D^{\prime}$. Пусть $\{B(y_i,r_i)\}$ ---


набор шаров, выбранный согласно лемме 2; $\{B(y_i,r_i)\}$ и


$\{B(y_i,2r_i)\}$ образуют покрытия множества $U$, и кратность


покрытия $\{B(y_i,2r_i)\}$ конечна ($B(y_i,2r_i)\subset U$ для


всех $i\in{\Bbb N}$). Тогда для функции $f_i$, ассоциированной с


шаром $B(y_i,r_i)$, имеем $\varphi^{\ast}f_i=1$ на множестве


$\varphi^{-1}(B(y_i,r_i))$ и $\varphi^{\ast}f=0$ вне множества


$\varphi^{-1}(B(y_i,2r_i))$, в частности, $\varphi^{\ast}f=0$ на


множестве $T$. Кроме того, для нее справедлива оценка


$$


\|\varphi^{\ast}f_i\mid L^1_q(2Q)\|\leq


\|\varphi^{\ast}f_i\mid L^1_q(D)\| \leq


C_1{\Phi(B(y_i,2r_i))}^{\frac{1}{\sigma}} |B(y_i,r_i)|^{\frac{1}{p} -


\frac{1}{n}}


$$


(при $q=p$ полагаем


$\Phi(A)^{\frac{1}{\sigma}}=\|\varphi^*\|$). В силу одного из


вариантов неравенства Пуанкаре (см., напр., [12]) справедливо


неравенство


$$


\bigg(\int_{Q}|g|^{q^*}\,dx\bigg)^{\frac{1}{{q^*}}}\leq


C_2l(Q)^{n/{q^*}} \bigg(\int_{2Q}|\nabla


g|^q\,dx\bigg)^{\frac{1}{q}},


$$


где ${q^*}=qn/(n-q)$ при $q<n$, $l(Q)$ --- длина ребра куба $Q$,


а $g\in W^1_{q,\loc}(Q)$ ---


произвольная функция, равная нулю на множестве $T$ (здесь


$C_2$ --- постоянная, зависящая от ``размера'' множества


$T$). Применяя неравенство Пуанкаре к функции $\varphi^{\ast}f_i$


вместо $g$ и учитывая две последние оценки, получим


$$


|\varphi^{-1}(B(y_i,r_i))\cap Q|^{\frac{1}{q} - \frac{1}{n}}\leq


C_3\Phi(B(y_i,2r_i))^{\frac{1}{\sigma}}


|B(y_i,r_i)|^{\frac{1}{p} - \frac{1}{n}}.


$$


С помощью неравенства Г\"ельдера выводим соотношение


$$


\bigg(\,


\sum_{i=1}^{\infty}|\varphi^{-1}(B(y_i,r_i))\cap Q|


\bigg)^{\frac{1}{q} - \frac{1}{n}}


\leq C_3\bigg(\,\sum_{i=1}^{\infty}\Phi(B(y_i,2r_i))


\bigg)^{\frac{1}{r \sigma}}


\bigg(\,\sum_{i=1}^{\infty}|B(y_i,r_i)|


\bigg)^{\frac{1}{p} - \frac{1}{n}}.


$$


Применяя лемму 3 для оценки значений функции $\Phi$ через


кратность покрытия, имеем


$$


|\varphi^{-1}(E)\cap Q|^{\frac{1}{q} -\frac{1}{n}}\leq


C_4\Phi(U)^{\frac{1}{\sigma}}


|U|^{\frac{1}{p} - \frac{1}{n}}.


$$


Из этой оценки (используя следствие 1 при $p=n$) в силу произвола в выборе


открытого множества $U$ получим, что $|\varphi^{-1}(E)\cap Q|=0$.


Так как выбор куба $Q\subset D$ произволен, то


$|\varphi^{-1}(E)|=0$, и поэтому


отображение $\varphi$ обладает ${\cal N}^{-1}$-свойством.


\end{pf}





Из теоремы 4 и формулы замены переменной получим





\begin{cor}


Пусть отображение $\varphi :


D\to D^{\prime}$ порождает ограниченный оператор вложения


$$


\varphi^{\ast} : L^1_p(D^{\prime})\to


L^1_q(D),\quad 1\leq q\leq p\leq n.


$$


Тогда его якобиан $J(x,\varphi)$ почти всюду отличен от нуля.


\end{cor}





\begin{pf}


Согласно теореме 4 отображение


$\varphi$ обладает ${\cal N}^{-1}$-свойством. Это эквивалентно


тому, что образ множества


$A\subset D$, $A\cap \Sigma_{\varphi}\ne\emptyset$,


положительной меры есть множество положительной меры.


Пусть


$Z=\{x\in D\setminus \Sigma_{\varphi}\mid J(x,\varphi)=0\}$.


Если $|Z|>0$, то $|\varphi(Z\setminus \Sigma_{\varphi})|>0$.


Однако по формуле (1)


$|\varphi(Z\setminus \Sigma_{\varphi})|=0$.


Полученное противоречие приводит к тому, что $|Z|=0$.


\end{pf}





В следующем утверждении описывается достаточно широкий класс отображений,


порождающих отображения классов Соболева. Напомним, что


$$


K_p(x)=\inf\{k:


|D\varphi|(x)\leq k|J(x,\varphi)|^{1/p},\ x\in D \}.


$$





\begin{thm}


Пусть непостоянное на каждой компоненте связности


открытого множества отображение 


$\varphi : D\to D^{\prime}$, где $D$ и $D^{\prime}$ --- открытые множества в


${\Bbb R}^n$,


\begin{itemize}


\item[a)] принадлежит классу $\ACL(D);$


\item[b)] имеет конечное искажение\/$;$


\item[c)] функция кратности


$M(y,\varphi,D)=\#\{x\in\varphi^{-1}(y)


\setminus\Sigma_\varphi\}\in L_{\infty}(D^{\prime});$


\item[d)] величина


$K_{p,q}(D)=\|K_p(\cdot)\mid L_{\varkappa}(D)\|$


конечна, где $\varkappa^{-1}= 1/q - 1/p$ при $q<p$ и


$\varkappa=\infty$ при $q=p$.


\end{itemize}


Тогда


\begin{itemize}


\item[1)] отображение $\varphi$ порождает ограниченный оператор


$$


\varphi^{*} :L_{p}^{1}(D^{\prime})\to L_{q}^{1}(D),\quad 1\leq


q\leq p\leq \infty,


$$


причем его норма не превосходит величины


$$


\|M(y,\varphi,D)\mid L_{\infty}(D')\|^{\frac{1}{p}}


\|K_p(\cdot)\mid L_{\varkappa}(D)\bigr\|,\quad \varkappa\in[1,\infty];


$$


\item[2)] при $1\leq q\leq p\leq n$ отображение $\varphi$ обладает


${\cal N}^{-1}$-свойством Лузина и якобиан $J(x,\varphi)$


почти всюду отличен от нуля.


\end{itemize}





При $q=p=n$ справедливо обращение\/$:$ если отображение $\varphi :


D\to D^{\prime}$ порождает ограниченный оператор $\varphi^{*}


:L_{n}^{1}(D^{\prime})\to L_{n}^{1}(D)$, то $\varphi$ обладает


свойствами $\mathrm{a)}$--$\mathrm{c)}$ и величина


$K_{n,n}(D)=\|K_n(\cdot)\mid L_{\infty}(D)\|$


конечна. Для нормы оператора


$\varphi^{*}:L_{n}^{1}(D^{\prime})\to L_{n}^{1}(D)$


справедлива оценка снизу


$$


\|\varphi^*\|\sim\|H_n(\cdot)\mid L_{\infty}(D)\bigr\|\geq


\|M(y,\varphi,D)\mid L_{\infty}(D^{\prime})\|^{\frac{1}{n}}.


$$


\end{thm}





\begin{pf}


Проверим, что выполнены условия теоремы 1.


Для этого надо установить ограниченность величины


$$


H_{p,q}(D^{\prime})=\|H_q(\cdot)\mid L_{\varkappa}(D')\|,


$$


где функция $y\mapsto H_q(y)$


определена формулой (2), а число $\varkappa$


определено в условии теоремы.


Приводимые ниже оценки соответствуют случаю $q<p$.


При $q=p$ они лишь упрощаются.








Применяя формулу (1), получим оценку


\begin{multline*}


\|K_p(\cdot)\mid L_{\varkappa}(D)\|^\varkappa= \int_{D}


\bigg(\frac{|D\varphi|^p(x)}{|J(x,\varphi)|}\bigg)^{\frac{q}{p-q}}\,dx= \\


%


=\int_{D\setminus Z}


\bigg(\frac{|D\varphi|^q(x)}{|J(x,\varphi)|}


\bigg)^{\frac{p}{p-q}}|J(x,\varphi)|\,dx=


\int_{D^{\prime}}\; \sum_{x\in\varphi^{-1}(y)\setminus


\Sigma_{\varphi}} \bigg(\frac{|D\varphi|^q(x)}{|J(x,\varphi)|}


\bigg)^{\frac{p}{p-q}} \,dy,


\end{multline*}


где обозначения $\Sigma_{\varphi}$ и $Z$ имеют тот же смысл, что и


ранее.





С другой стороны, очевидно, что для почти всех $y\in D^{\prime}$ имеем


$$


\bigg(\,\sum_{x\in\varphi^{-1}(y)\setminus


\Sigma_{\varphi}} \frac{|D\varphi|^q(x)}{|J(x,\varphi)|}


\bigg)^{\frac{p}{p-q}} \leq M(y,\varphi,D)^{\frac{\varkappa}{p}}


\sum_{x\in\varphi^{-1}(y)\setminus\Sigma_{\varphi}}


\bigg(\frac{|D\varphi|^q(x)}{|J(x,\varphi)|}\bigg)^{\frac{p}{p-q}}.


$$


Из приведенных оценок непосредственно следует


$$


\|H_q(\cdot)\mid L_{\varkappa}(D')\|\leq


\|M(\cdot,\varphi,D)\mid L_{\infty}(D')\|^{\frac{1}{p}}


\|K_p(\cdot)\mid L_{\varkappa}(D)\|,\quad \varkappa\in[1,\infty].


$$


Отсюда видно, что


выполнены условия теоремы 1, а следовательно, и теоремы 4. Таким


образом, из теоремы 1 получаем первое заключение теоремы 5, а из


теоремы 4 и следствия 4 вытекают ${\cal N}^{-1}$-свойство и


невырожденность якобиана почти всюду.





Если теперь $q=p=n$ и отображение $\varphi : D\to D^{\prime}$


порождает ограниченный оператор $\varphi^{*}


:L_{n}^{1}(D^{\prime})\to L_{n}^{1}(D)$, то из теоремы 1 вытекают


свойства a), b) и


$H_{n,n}(D^{\prime})=\|H_n(\cdot)\mid


L_{\infty}(D^{\prime})\|<\infty$. Отсюда получим


$K_{n,n}(D)=\|K_n(\cdot)\mid L_{\infty}(D)\|<\infty$,


т.\,к. $K_{n,n}(D)\leq H_{n,n}(D^{\prime})$ в силу поточечного


неравенства $K_n(x)\leq H_n(\varphi(x))$, справедливого для почти


всех $x\in D\setminus\Sigma_\varphi$, поскольку по следствию 4


якобиан $J(x,\varphi)$ почти всюду отличен от нуля.





Чтобы доказать свойство


c), заметим, что $|J(x,\varphi)|\leq|D\varphi|^n(x)$


почти всюду, и поэтому


$$


M(y,\varphi,D)^{\frac{1}{n}}\leq H_n(y)\leq


\|H_n(\cdot)\mid L_{\infty}(D^{\prime})\|


$$


для почти всех $y\in D^{\prime}$.


\end{pf}





В том случае, когда отображение $\varphi$ удовлетворяет


дополнительным топологическим требованиям, можно


получить точную оценку искажения меры.





Напомним, что отображение $\varphi : D\to D^{\prime}$


называется {\it собственным}, если прообраз компактного


множества из $D^{\prime}$ компактен в $D$.





\begin{thm}


Пусть отображение $\varphi :


D\to D^{\prime}$ порождает ограниченный оператор вложения


$$


\varphi^{\ast} : L^1_p(D^{\prime})\to L^1_q(D),\quad 1\leq q\leq


p\leq n.


$$


Тогда при $p<n$ для всякого измеримого множества


$E\subset D^{\prime}$ выполняется неравенство


$$


|\varphi^{-1}(E)|^{\frac{1}{q} - \frac{1}{n}} \leq \|H_q(\cdot)\mid


L_\varkappa(E)\| |E|^{\frac{1}{p} -\frac{1}{n}}.


$$


Если отображение


$\varphi : D\to D^{\prime}$ собственное и порождает ограниченный


оператор вложения $\varphi^{\ast} : L^1_n(D^{\prime})\to


L^1_q(D)$, $q<n$, то


$$


|\varphi^{-1}(E)|^{\frac{1}{q} -\frac{1}{n}}


\leq \|H_q(\cdot)\mid L_\varkappa(E)\|


$$


для любого измеримого


множества $E\subset D^{\prime}$. Здесь $\varkappa=\frac{pq}{p-q}$,


если $q<p$, и $\varkappa=\infty$, если $q=p$.


\end{thm}





{\bf Доказательство.}


В случае $p<n$ воспользуемся п.\,2


теоремы 3, полагая в ней $r=\frac{np}{n-p}$, a $A^{\prime}=E$.


Тогда $s=\frac{nq}{n-q}$


и для функции $f=\chi_E$ в силу ограниченности оператора


$\varphi^{\ast} : L_r(E)\to L_s(\varphi^{-1}(E))$ имеем


$|\varphi^{-1}(E)|^{\frac{1}{q} - \frac{1}{n}}


\leq \|H_q(\cdot)\mid L_\varkappa(E)\|\,|E|^{\frac{1}{p} - \frac{1}{n}}$.





В случае $q<p=n$ приведем независимое доказательство. Фиксируем


срезку $\eta\in C_{0}^{\infty}({\Bbb R}^n)$, равную единице на


$B(0,1)$ и нулю вне шара $B(0,2)$. На основании леммы~1 имеем


($B=B(y_0,r)$, $\ov{B(y_0,2r)}\subset D'$)


$$


\|\varphi^{\ast}f\mid L^1_q(D)\| \leq


{\Phi(2B)}^{\frac{1}{\varkappa}},


$$


где $f(y)=\eta\big(\frac{y-y_0}{r}\big)$. Так как отображение


$\varphi$ собственное, то


$\varphi^{\ast}f\in\overset{\circ}{L}{}^1_q(D)$ и


выполнено неравенство Соболева


$$


\|\varphi^{\ast}f\mid


L_{q^{\ast}}(D)\| \leq C_1 \|\varphi^{\ast}f\mid


L^1_q(D)\|,\quad\text{где}\quad


\frac{1}{q^{\ast}}=\frac{1}{q} -\frac{1}{n}


$$


(см., напр., [12]). Из условий $\varphi^{\ast}f=1$ на множестве


$\varphi^{-1}(B(y_0,r))$ и $\varphi^{\ast}f=0$ вне множества


$\varphi^{-1}(B(y_0,2r))$ получим оценку


$$


|\varphi^{-1}(B)|^{\frac{1}{q} - \frac{1}{n}}\leq


C_2\Phi(2B)^{\frac{1}{\varkappa}} |B|.


$$


Для измеримого множества


$E\subset D^{\prime}$ выберем число $\varepsilon>0$ и


открытое множество $U\supset E$ так, чтобы


$|U|\leq|E|+\varepsilon$. Рассмотрим набор шаров $\{B(y_i,r_i)\}$


такой, что $\{B(y_i,r_i)\}$ и $\{B(y_i,2r_i)\}$ образуют покрытия


множества $U$, кратность которых конечна и не зависит от выбора


открытого множества $U$ ($B(y_i,2r_i)\subset U$ для всех $i\in\Bbb N$). С


помощью неравенства Г\"ельдера выводим


$$


\bigg(\,\sum_{i=1}^{\infty}|\varphi^{-1}(B(y_i,r_i))|


\bigg)^{\frac{1}{q} - \frac{1}{n}} \leq C_2


\bigg(\,\sum_{i=1}^{\infty}


\int_{B(y_i,2r_i)}H_q(y)^{\varkappa}\,dy


\bigg)^{\frac{1}{\varkappa}}.


$$


Отсюда имеем


$|\varphi^{-1}(E)|^{\frac{1}{q} - \frac{1}{n}} \leq


C\|H_q(\cdot)\mid L_{\varkappa}(U)\|$ в силу теоремы 1. Так как


число $\varepsilon>0$ было выбрано произвольно, то, переходя к


пределу при $\varepsilon \to 0$, получим


$$


|\varphi^{-1}(E)|^{\frac{1}{q} - \frac{1}{n}} \leq


C\|H_q(\cdot)\mid L_{\varkappa}(E)\|.\quad\square


$$





\begin{notenonum}


Требование для отображения $\varphi:D\to D^{\prime}$


быть собственным в теореме 6 при $p=n$


вызвано не только применением теоремы вложения, но и существом дела


(см. [32]).


Естественно, что оценка теоремы 6 без топологического условия на


отображение справедлива также и в тех случаях, когда на множестве


определения отображения справедливы такие же теоремы вложения, как


и в пространстве ${\Bbb R}^n$.


\end{notenonum}








\section{Топологические отображения с интегрируемым искажением}





В этом пункте покажем, что при гомеоморфных отображениях норма


оператора вложения ассоциируется со счетно-квазиаддитивной


функцией множества, определенной на открытом множестве $D$.





Для открытого множества $A\subset D\subset{\Bbb R}^n$ обозначим


символом $\varphi_A$ сужение гомеоморфизма $\varphi :


D\to D^{\prime}$ на множество $A$, и символом $\varphi_A^{*}$


оператор суперпозиции $f\circ\varphi$, где $f\in


L^1_p(\varphi(A))$.





Определим функцию $\Psi(A)$ открытого множества $A\subset D$ по


правилу


$$


\Psi(A)=\sup_{f\in {L}_{p}^{1}(\varphi(A))}


\left( \frac{\|\varphi^{\ast} f\mid {L}_{q}^{1}(A)\|}


{\|f\mid {L}_{p}^{1}(\varphi(A))\|}


\right)^{\sigma},\quad \text{где}\quad \sigma=


\begin{cases}


\frac{pq}{p-q} & \text{при}\ 1\leq q<p<\infty;\\


q & \text{при}\ 1\leq q< p=\infty;\\


1 & \text{при}\ 1\leq q =p<\infty.


\end{cases}


$$


Таким образом, по определению


$\Psi(A)^{\frac{1}{\sigma}}=\|\varphi^*_A\|$. Корректность определения


функции $\Psi$ следует из работ [27] и [33], в которых доказана


оценка


$$


\Psi(A)\leq


\begin{cases} \int\limits_A


K_p^\sigma(x)\,dx & \text{при}\ 1\leq q<p<\infty; \\


\esssup\limits_{x\in A} K_p(x) & \text{при}\ 1\leq q=p<\infty;\\


\int\limits_A |D\varphi|^q(x)\,dx &


\text{при}\ 1\leq q<p=\infty,


\end{cases}


$$


где величина $K_p(x)$ определена перед теоремой 5.


Из этих неравенств получаем, что из ограниченности оператора


вложения


$$


\varphi^{*} :L_{p}^{1}(D^{\prime})\to


L_{q}^{1}(D),\ \ 1\leq q\leq p\leq \infty,


$$


вытекает ограниченность оператора вложения


$$


\varphi^{*}_A


:L_{p}^{1}(\varphi(A))\to L_{q}^{1}(A),\ \ 1\leq q\leq p\leq


\infty.


$$


Следовательно, при $q<p$ функция множества $\Psi$


является ограниченной абсолютно непрерывной функцией, определенной


на открытых подмножествах из $D$.








\begin{lem}


Пусть гомеоморфизм $\varphi :D\to D^{\prime}$


порождает ограниченный оператор


вложения $\varphi^{*} :L_{p}^{1}(D^{\prime})\to L_{q}^{1}(D)$,


$1\leq q<p\leq \infty$.


Тогда функция $\Psi$ является ограниченной абсолютно непрерывной


счетно-аддитивной функцией,


определенной на открытых множествах из~$D$.


\end{lem}





\begin{pf}


Пусть $A_{i}$, $i\in {\Bbb N}$, ---


открытые попарно непересекающиеся множества в $D$,


$A_i^{\prime}=\varphi(A_{i})$, $i=0,1,\dotsc$,


$A^{\prime}_0=\bigcup\limits_{i=1}^{\infty}A_i^{\prime}$.


Рассмотрим такую функцию $f_i\in {L}_{p}^{1}


(A_i^{\prime})$, чтобы одновременно выполнялись условия


$\|\varphi^{\ast} f_i\mid {L}_{q}^{1}


(A_i)\|\geq \big(\Psi(A_i^{\prime})


(1-\frac{\varepsilon}{2^i})\big)^{\frac{1}{\sigma}}


\|f_i\mid {L}_{p}^{1}(A_i^{\prime})\|$


и


$\|f_i\mid {L}_{p}^{1}(A_i^{\prime})


\|^p=\Psi(A_i^{\prime})\big(1-\frac{\varepsilon}{2^i}\big)$


при $p<\infty$


($\|f_i\mid {L}_{p}^{1}(A_i^{\prime})


\|=1$ при $p=\infty$), $\varepsilon\in(0,1)$.


Пусть $f_N=\sum\limits_{i=1}^{N}f_i$.


Повторяя далее рассуждения леммы 1, получим


$$


\Psi(A_0^{\prime})^{\frac{1}{\sigma}}\geq\sup\frac


{\bigg\|\varphi^{\ast} f_N\,\Big|\, {L}_{q}^{1}


\Big(\bigcup\limits_{i=1}^{N}A_i\Big)\bigg\|}


{\bigg\|f_N\,\Big|\, {L}_{p}^{1}


\Big(\bigcup\limits_{i=1}^{N}A_i^{\prime}\Big)\bigg\|}\geq


\bigg(\,\sum_{i=1}^{N}\Psi(A_i^{\prime})-\varepsilon\Psi(A_0^{\prime})


\bigg)^{\frac{1}{\sigma}},


$$


где точная верхняя грань берется по всем функциям


$f_N\in {L}_{p}^{1}


\Big(\bigcup\limits_{i=1}^{N}A_i^{\prime}\Big)$


указанного выше вида. Так как $N$ и $\varepsilon$ произвольны, то


$$


\sum_{i=1}^{\infty}\Psi(A^{\prime}_i) \leq


\Psi\Big(\bigcup_{i=1}^{\infty}A_i^{\prime}\Big).


$$


Справедливость обратного неравенства проверяется непосредственно.


\end{pf}





Следующая теорема для любого открытого множества $A\subset D$ 


уточняет оценку нормы оператора вложения


$$


\varphi^{\ast}_A : L^{1}_p(\varphi(A))\to L^1_q(A),\quad 1\leq q\leq


p\leq\infty,


$$


полученную в [27];


{\it функция множества $\Psi$ хоть и не является монотонной, тем не менее


заключена между двумя монотонными счетно-аддитивными функциями множества}.





Определим функцию $\Lambda(A)$ открытого множества $A\subset


D$ по правилу


$$


\Lambda(A)=\sup_{f\in\overset{\circ}{L}{}_{p}^{1}(\varphi(A))}


\left( \frac{\|\varphi^{\ast} f\mid {L}_{q}^{1}(A)\|}


{\|f\mid \overset{\circ}{L}{}_{p}^{1}(\varphi(A))\|}


\right)^{\sigma},\quad \text{где}\ \ \sigma=


\begin{cases}


\frac{pq}{p-q} & \text{при}\ 1\leq q<p<\infty;\\


1 & \text{при}\ 1\leq q = p<\infty.


\end{cases}


$$


Так же, как и в лемме 1,


доказывается, что $\Lambda$ является ограниченной монотонной


счетно-аддитивной функцией, определенной на открытых ограниченных


множествах $A\subset D$. Очевидно $\Lambda(A)\leq\Psi(A)$ и


$\Lambda(A)\leq\Psi(B)$ для любого ограниченного открытого


множества $A$ и произвольного открытого множества $B$ с условием


$A\subset B\subset D$.





\begin{thm}


Пусть $D$ и $D^{\prime}$ ---


открытые множества в ${\Bbb R}^n$. Гомеоморфизм $\varphi :D\to


D^{\prime}$ порождает ограниченный оператор вложения $\varphi^{*}


:L_{p}^{1}(D^{\prime})\to L_{q}^{1}(D)$, $1\leq q\leq p<\infty$,


тогда и только тогда, когда отображение $\varphi$ принадлежит


$W^1_{q,\loc}(D)$, имеет конечное искажение и


$$


K_p(\cdot)\in


L_\varkappa(D),\quad\text{где}\quad \varkappa=


\begin{cases}


\frac{qp}{p-q}, &\text{если}\ q<p;\\


\infty, &\text{если}\ q=p.


\end{cases}


$$





Оператор


$\varphi^{*}_A :L_{p}^{1}(\varphi(A))\to L_{q}^{1}(A)$, $A\subset D$ ---


открытое множество, ограничен и его норма


удовлетворяет соотношению


\begin{equation}


\gamma_{p,q} K_{p,q}(A)\leq\|\varphi^{*}_A \| \leq


K_{p,q}(A),


\end{equation}


где $K_{p,q}(A)=\|K_p(\cdot)\mid


L_\varkappa(D)\|$, а $0<\gamma_{p,q}$ --- постоянная.


\end{thm}





Для доказательства теоремы 7 при $q<p$ понадобится


представляющая самостоятельный интерес





\begin{thm}


Пусть $D$ и $D^{\prime}$ ---


открытые множества в ${\Bbb R}^n$. Гомеоморфизм $\varphi :D\to


D^{\prime}$ порождает ограниченный оператор вложения $\varphi^{*}


:L_{p}^{1}(D^{\prime})\to L_{q}^{1}(D)$, $1\leq q< p<\infty$,


тогда и только тогда, когда отображение $\varphi$ принадлежит


$W^1_{q,\loc}(D)$, имеет конечное искажение и существует


ограниченная монотонная счетно-аддитивная функция $\Lambda$,


определенная на открытых ограниченных множествах $A\subset


D$, такая, что справедлива эквивалентность


$$


\gamma^{\varkappa}_{p,q} K^{\varkappa}_{p}(x)\leq\Lambda^{\prime}(x)


\leq K^{\varkappa}_p(x),


$$


где $\gamma_{p,q}$ --- постоянная из


теоремы $7$.


\end{thm}





\begin{pf*}{Доказательство теорем 7 и 8. Необходимость}


Принадлежность отображения $\varphi$ классу


$W^1_{q,\loc}(D)$ проверяется подстановкой тестовых функций,


которые применялись при доказательстве предложения 2. Покажем, что


отображение $\varphi$ имеет конечное искажение. Заметим, что этот


факт так же, как и оценка коэффициента квазиконформности


отображения $\varphi$ через производную функции множества,


заданной на подмножествах открытого множества $D^{\prime}$, может быть


получен как следствие теоремы 2, доказанной при более общих


предположениях. Однако мы приведем здесь независимое


рассуждение.





Так же, как и в доказательстве теоремы 1, для любой функции


$f\in \overset{\circ}{L}{}_{p}^{1}(A')$ при $q\leq p$


выполняется неравенство


$$


\|\varphi^{\ast}f\mid{L}_{q}^{1}(D)\|


\leq\Lambda(\varphi^{-1}(A'))^{1/\sigma}


\|f\mid\overset{\circ}{L}{}_{p}^{1}(A')\|,


$$


где $A'\subset D^{\prime}$ --- открытое ограниченное


подмножество. Далее полагаем $q<p$ (при $q=p$ выкладки лишь


упрощаются). Фиксируем срезку $\eta\in C_{0}^{\infty}({\bf R}^n)$,


равную единице на $B(0,1)$ и нулю вне шара $B(0,2)$. Подставляя в


это неравенство функции


$h_{j}(z)=(z-y)_{j}\eta(\frac{z-y}{r})$, $j=1,\dotsc,n$, где


$(z-y)_{j}$ обозначает $j$-ю координату вектора $z-y$,


$B(y,2r)\subset D^{\prime}$, имеем


$$


\bigg(\,\int\limits_{\varphi^{-1}(B(y,r))}|D\varphi|^q(x)\,dx


\bigg)^{1/q}\leq


C_1\Lambda(\varphi^{-1}(B(y,2r)))^{1/\varkappa}|B(y,r)|^{1/p},


$$


где $C_1$ --- некоторая постоянная. Пусть $B(x_0,r)\subset D$ и


$U=\varphi(B(x_0,r))$. Рассмотрим набор шаров $\{B(y_i,r_i)\}$


такой, что $\{B(y_i,r_i)\}$ и $\{B(y_i,2r_i)\}$ образуют покрытие


множества $U$, кратность которых конечна и не зависит от


выбора открытого множества $U$ ($B(y_i,2r_i)\subset U$ для всех


$i\in{\Bbb N}$). Применяя предыдущее неравенство к каждому из


шаров $\{B(y_i,2r_i)\}$, c помощью неравенства Г\"ельдера выводим


соотношение


$$


\bigg( \sum_{i=1}^{\infty}


\int\limits_{\varphi^{-1}(B(y_i,r_i))}|D\varphi|^q(x)\,dx


\bigg)^{\frac{1}{q}} \leq C_2


\bigg(\sum_{i=1}^{\infty}\Lambda(\varphi^{-1}


(B(y_i,2r_i)))\bigg)^{\frac{1}{\varkappa}}


\bigg(\sum\limits_{i=1}^{\infty}|B(y_i,r_i)|


\bigg)^{\frac{1}{p}}.


$$


Заметим, что функция множества


$A'\mapsto\Lambda(\varphi^{-1}(A'))$, $A'\subset D'$~--- открытое


множество, является монотонной и счетно-аддитивной. Применяя к ней


лемму 3, получим оценку для суммы значений $\Lambda(\varphi^{-1}


(B(y_i,2r_i)))$ через функцию кратности


$$


\sum_{i=1}^{\infty}\Lambda(\varphi^{-1}(B(y_i,2r_i)))\leq


\xi_n\Lambda\Big(\varphi^{-1}\Big(\bigcup_{i=1}^{\infty}


B(y_i,2r_i)\Big)\Big),


$$


где постоянная $\xi_n$ зависит только от размерности $n$ [31].


Учитывая, что


$U=\varphi(B(x_0,r))=\bigcup\limits_{i=1}^{\infty}B(y_i,2r_i)=


\bigcup\limits_{i=1}^{\infty}B(y_i,r_i)$,


имеем


$$


\bigg(\, \int\limits_{B(x_0,r)}|D\varphi|^q(x)\,dx


\bigg)^{\frac{1}{q}} \leq C_3 \Lambda(B(x_0,r))^{\frac{1}{\varkappa}}


|\varphi(B(x_0,r))|^{\frac{1}{p}}.


$$


Преобразуем это неравенство к виду


$$


\bigg(\frac{1}{|B(x_0,r)|}\int\limits_{B(x_0,r)}|D\varphi|^q(x)\,dx


\bigg)^{1/q}\leq


C_3\bigg({\frac{\Lambda(B(x_0,r))}{|B(x_0,r)|}}\bigg)^{1/\varkappa}


\bigg({\frac{|\varphi(B(x_0,r))|}{|B(x_0,r)|}}\bigg)^{1/p}


$$


и перейдем к пределу при $r\to 0$. Тогда


$$


|D\varphi|(x)\leq C_3


\begin{cases}


\Lambda^{\prime}(x)^{1/\varkappa}


J_{\varphi}(x)^{1/p} & \text{при}\ q<p;\\


\|\varphi^*_D\| J_{\varphi}(x)^{1/p} & \text{при}\ q=p.


\end{cases}


$$





Так как отображение $\varphi$ принадлежит $W^1_{1,\loc}(D)$, то


(напр., [29], [34]) существует последовательность


замкнутых множеств $\{A_k\}$, $A_k\subset A_{k-1}$, таких, что


$\varphi|_{A_k}\in\Lip (A_k)$ и $|D\setminus \bigcup\limits_k


A_k|=0$. Отсюда получим, в частности, что


$J_{\varphi}(x)=J(x,\varphi)$ для почти всех $x\in D$


([29]), и, следовательно, отображение $\varphi$ имеет конечное


искажение. Кроме того,


\begin{alignat}{2}


K_p(x)&=\frac{|D\varphi|(x)}{|J(x,\varphi)|^{1/p}}\leq


c\|\varphi^{\ast}_D\| &\ \ &\text{для почти всех}\ x\in D \ \text{при}


\ 1\leq q=p<\infty \\


%


\intertext{и}


%


K^\varkappa_p(x)&=\bigg(\frac{|D\varphi|(x)}{|J(x,\varphi)|^{1/p}}


\bigg)^\varkappa


\leq c\Lambda^{\prime}(x) &\ \ &\text{для почти всех}\ x\in D\ \text{при}


\ 1\leq q<p<\infty.


\end{alignat}


Интегрируя (8) по ограниченному


открытому множеству $A\subset D$ и применяя


предложение 1, получим оценку нормы оператора


$\varphi_D^{\ast}$ снизу


\begin{equation}


\int_A\bigg(\frac{|D\varphi|(x)}{|J(x,\varphi)|^{1/p}}\bigg)^\varkappa\,dx


\leq c\int_A\Lambda^{\prime}(x)\,dx\leq c\Lambda(A)\leq


c\Psi(D).


\end{equation}





{\bf Достаточность}. Пусть $A\subset D$ --- произвольное открытое


множество, а $A^{\prime}=\varphi(A)$. Покажем, что для любой функции


$f\in L_p^1(A^{\prime})\cap C^1(A^{\prime})$ выполняется неравенство


$\|\varphi^*f\mid L_q^1(A)\|\le K_{p,q}\|f\mid L_p^1(A^{\prime})\|$,


$q\leq p$. Так как $f\circ\varphi$ принадлежит классу $\ACL(A)$, то


\begin{multline*}


\|\varphi^* f\mid L_q^1(A)\|\le


\bigg(\int_A(|\nabla


f|(\varphi(x))|D\varphi|)^q\,dx\bigg)^{1/q}= \\


%


= \bigg(\int_{A\setminus Z}|\nabla


f|^q(\varphi(x))|J(x,\varphi)|^{q/p}


\frac{|D\varphi|^q}{|J(x,\varphi)|^{q/p}}\,dx\bigg)^{1/q}.


\end{multline*}


Используя неравенство Г\"ельдера при $q\leq p$, выводим оценку


$$


\|\varphi^*f\mid L_q^1(A)\|\le


\bigg(\int_{A\setminus Z}\bigg(\frac{|D\varphi|(x)}


{|J(x,\varphi)|^{1/p}}\bigg)^{\varkappa}\,dx\bigg)^{1/\varkappa}


\bigg(\int_{A\setminus Z}|\nabla f|^p(\varphi(x))


|J(x,\varphi)|\,dx\bigg)^{1/p}


$$


(при $q=p$ левый сомножитель


равен $K_{p,p}(A)$). Применяя к правому сомножителю формулу (1),


получим оценку для нормы $\|\varphi^{\ast}_A\|\leq K_{p,q}(A)$.


Распространение полученной оценки на все функции $f\in


L_p^1(A^{\prime})$ происходит по схеме теоремы 1.





Соотношение $\Psi^{\frac{1}{\varkappa}}(A)=


\|\varphi^{\ast}_A\|\leq K_{p,q}(A)$ при $q\leq p$,


таким образом, доказано. Отсюда в силу неравенства


$\Lambda(A)\leq\Psi(A)$ получим оценку сверху для


производной $\Lambda^{\prime}(x)$ при $q<p$


$$


\Lambda^{\prime}(x)\leq


\bigg(\frac{|D\varphi|(x)}{|J(x,\varphi)|^{1/p}}\bigg)^\varkappa


\quad\text{для почти всех}\ x\in D\  \text{в случае}\ 1\leq


q<p<\infty.


$$





%% Докажем теперь левую часть соотношений (6).


Поскольку оператор $\varphi^{*}_A :L_{p}^{1}(\varphi(A))\to


L_{q}^{1}(A)$ ограничен, то из рассуждений первой части


доказательства, примененной к открытому множеству $A$ вместо $D$,


получаем для него оценку снизу; из (7) имеем


$$


\esssup_{x\in A}


\frac{|D\varphi|(x)}{|J(x,\varphi)|^{1/p}} \leq


C\|\varphi^{\ast}_A\|\quad \text{в случае}\ 1\leq q=p<\infty,


$$


а в силу (9) и неравенства $\Lambda(A)\leq \Psi(A)$ ---


$$


\int_A\bigg(\frac{|D\varphi|(x)}{|J(x,\varphi)|^{1/p}}\bigg)^\varkappa\,dx


\leq c\Psi(A)


$$


для любого ограниченного множества $A\subset D$.


Если $B\subset D$ --- произвольное открытое


множество, то для


произвольной возрастающей последовательности $A_i\subset


A_{i+1}\subset B$ ограниченных открытых множеств таких, что


$\bigcup\limits_{i=1}^\infty A_i=B$, имеем


$K^\varkappa_{p,q}(A_i)\leq c\Lambda(A_i)\leq c\Psi(B)$. Переходя к


пределу при $i\to\infty$, получаем необходимую оценку.


\end{pf*}





\begin{cor}


Функция множества $\Psi$


является ограниченной абсолютно непрерывной счетно


$\vartheta$-квазиаддитивной функцией множества, определенной на


открытых множествах из $D$.


\end{cor}





{\bf Доказательство.}


%% Так как функция $\Psi$ ограничена, то по


%% теореме 7 отображение $\varphi$ имеет конечное искажение.


Для произвольной совокупности попарно непересекающихся


открытых множеств $A_i\subset A\subset D$ имеем


$$


\sum_{i=1}^\infty\Psi(A_i)\leq


\sum_{i=1}^\infty\int_{A_i}


\bigg(\frac{|D\varphi|(x)}{|J(x,\varphi)|^{1/p}}\bigg)^\varkappa\,dx\leq


\int_{A}\bigg(\frac{|D\varphi|(x)}{|J(x,\varphi)|^{1/p}}\bigg)^\varkappa\,dx


\leq \vartheta\Psi(A).


$$





Отсюда, из теорем 7 и 8 и предложения 1 получим





\begin{cor}


Пусть $D$ и $D^{\prime}$ ---


открытые множества в ${\Bbb R}^n$. Гомеоморфизм $\varphi :D\to D^{\prime}$


порождает ограниченный оператор вложения $\varphi^{*}


:L_{p}^{1}(D^{\prime})\to L_{q}^{1}(D)$, $1\leq q<p<\infty$, тогда


и только тогда, когда существует ограниченная абсолютно


непрерывная счетно $\vartheta$-квазиаддитивная функция $\Psi$,


определенная на открытых множествах $A\subset D$, причем ее


верхняя производная $\ov\Psi{}^{\prime}(x)$ эквивалентна


$K^\varkappa_p(x)=


\Big(\frac{|D\varphi(x)|}{|J(x,\varphi)|^{1/p}}\Big)^\varkappa{:}$


$$


\lambda_2 K^\varkappa_p(x)\leq \ov\Psi{}^{\prime}(x)\leq


K^\varkappa_p(x)


$$


почти всюду в $D$, где $\lambda_2$ --- некоторая постоянная.





Норма оператора $\varphi^{*}_A$ эквивалентна


$\int\limits_A\ov\Psi{}^{\prime}(x)\,dx$


для любого открытого множества $A\subset D$.


\end{cor}





\begin{thebibliography}{99}





\bibitem{1}


Водопьянов С.К., Гольдштейн В.М. {\it Структурные


изоморфизмы пространств $W^1_n$ и квазиконформные отображения} //


Сиб. матем. журн. -- 1975. -- Т.\,16. -- \No\,2. -- С.\,224--246.





\bibitem{2}


Данфорд Н., Шварц Дж.Т.


{\it Линейные операторы. Общая теория.}


-- М.: Ин. лит., 1962. -- 895~с.





\bibitem{3}


Nakai M. {\it Algebraic criterion on quasiconformal equivalence of 


Riemann surfaces} //


Nagoya Math. J. -- 1960. -- V.\,16. -- P.\,157--184.





\bibitem{4}


Lewis L.G.


{\it Quasiconformal mappings and Royden algebras in space} //


Trans. Amer. Math. Soc. -- 1971. -- V.\,158. -- \No\,2. -- P.\,481--492.





\bibitem{5}


Водопьянов С.К., Гольдштейн В.М. {\it Функциональные


характеристики квазиизометрических отображений} // Сиб.


матем. журн. -- 1976. -- Т.\,17. -- \No\,4. -- С.\,768--773.





\bibitem{6}


Водопьянов С.К., Гольдштейн В.М. {\it Новый функциональный


инвариант для квазиконформных отображений} // Некоторые вопросы


современной теории функций: Материалы конф. -- Новосибирск. --


1976. -- С.\,18--20.





\bibitem{7}


Водопьянов С.К. {\it Отображения однородных групп и вложения


функциональных пpостранств} // Сиб. матем. журн. -- 1989. --


Т.\,30. -- \No\,5. -- С.\,25--41.





\bibitem{8}


Гольдштейн В.М., Романов А.С. {\it Об отображениях,


сохраняющих пространства Соболева} // Сиб. матем. журн. -- 1984.


-- Т.\,25. -- \No\,3. -- С.\,55--61.





\bibitem{9}


Романов А.С.


{\it О замене переменной в пространствах потенциалов Бесселя и Рисса} //


Функц. анал. и мат. физика.


Ин-т матем. СО АН СССР. -- Новосибирск.


-- 1985. -- С.\,117--133.





\bibitem{10}


Водопьянов С.К.


{\it $L_p$-теория потенциала и квазиконформные отображения на 


однородных группах} /  С.С.\,Кутателадзе.


Соврем. пробл. геометрии и анализа. --


Новосибирск: Наука, 1989. -- С.\,45--89.





\bibitem{11}


Мазья В.Г., Шапошникова Т.О. {\it Мультипликаторы в


пространствах дифференцируемых функций.} -- Л.:


Изд-во Ленингр. ун-та, 1986. -- 403~с.





\bibitem{12}


Мазья В.Г. {\it Пространства С.Л.\,Соболева.}


-- Л.: Изд-во Ленингр. ун-та, 1985. -- 416~с.





\bibitem{13}


Reiman~M.


{\it \"Uber harmonische Kapazit\"at und quasikonforme Abbildungen in Raum}


// Comm. Math. Helv. -- 1969. -- V.\,44. -- P.\,264--307.





\bibitem{14}


Gehring F.W.


{\it Lipschitz mappings and the $p$-capacity of rings in $n$-space} //


Advances in the Theory of Riemann surfaces, Ann. Math. Stud. -- 1971.


-- N.\,66. -- P.\,175--193. %%%%%% N. ??





\bibitem{15}


Lelong-Ferrand~J.


{\it Etude d'une classe d'applications li\'{e}es \`{a}


des homomorphismes d'alg\`{e}bres de fonctions et g\'{e}n\'{e}ralisant les


quasi-conformes} //


Duke Math. J. -- 1973. -- V.\,40. -- P.\,163--186.





\bibitem{16}


Водопьянов С.К. {\it Формула Тейлора и функциональные


пространства.} -- Новосибирск: Изд-во Новосибирcк. ун-та, 1988. -- 96~с.





\bibitem{17}


Водопьянов С.К.


{\it Весовые пространства Соболева и теория отображений} //


Тез. докл. Всесоюз. матем. шк. ``Теория потенциала''.


Кацивели, 26 июня--3 июля 1991.


-- Киев: Ин-т математики АН УССР. -- 1991. -- С.\,7.





\bibitem{18}


Водопьянов С.К. {\it Геометрические свойства пространств


функций c обобщенной


гладкостью}. Дис. $\dotsc$ докт. физ.-матем. наук.


-- Новосибирск: Ин-т математики им. С.Л.\,Со\-болева, 1992. -- 335~с.





\bibitem{19}


Mostow G.D. {\it Quasiconformal mappings in $n$-spaces 


and the rigidity of hyperbolic spaces forms} //


Inst. hautes \'etudes scient. publications mathematiques.


-- 1968. -- V.\,34. -- P.\,53--104.





\bibitem{20}


V\"ais\"al\"a~J. {\it Lectures on $n$-dimensional quasiconformal


mappings}. -- Berlin etc.: Springer. (Lecture Notes in Math.;


229), 1971. -- 144~p.





\bibitem{21}


Решетняк Ю.Г. {\it Пространственные отображения с


ограниченным искажением.} -- Новосибирск: Наука, 1982. --


279~с.





\bibitem{22}


Gol'dshte\u\i{}n V., Gurov L., Romanov A.


{\it Homeomorphisms that induce monomorphisms of Sobolev spaces} //


Israel J. Math. -- 1995. -- V.\,91. -- \No\,1--3. -- P.\,31--60.





\bibitem{23}


Ухлов А.Д. {\it Отображения, порождающие вложения пространств


Соболева} // Сиб. матем. журн. -- 1993. -- Т.\,34. -- \No\,1. --


С.\,185--192.





\bibitem{24}


Кругликов В.И. {\it Емкости конденсаторов и пространственные 


отображения, квазиконформные в среднем} //


Матем. сб. -- 1986. -- Т.\,130. -- \No\,2. -- С.\,185--206.





\bibitem{25}


Gol'dshte\u\i{}n V., Gurov L. {\it Applications of change of variables 


operators for exact embedding theorems} // Integral equations operator 


theory. -- 1994. -- V.\,19. -- \No\,1. -- P.\,1--24.





\bibitem{26}


Водопьянов С.К. {\it Топологические и геометрические свойства 


отображений классов Соболева с суммируемым якобианом}. I //


Сиб. матем. журн. -- 2000. -- Т.\,41. -- \No\,1.-- С.\,23--48.





\bibitem{27}


Водопьянов С.К., Ухлов А.Д. {\it Пространства Соболева и


$(P,Q)$-квазиконформные отображения групп Карно} // Сиб.


матем. журн. -- 1998. -- Т.\,39. -- \No\,4. -- С.\,776--795.





\bibitem{28}


Водопьянов С.К., Ухлов А.Д. {\it Операторы суперпозиции в


пространствах Лебега и дифференцируемость квазиаддитивных функций


множества} // Владикавказск. матем. журн. -- 2002. -- Т.\,4. --


\No\,1. -- С.\,11--33.





\bibitem{29}


Vodop'yanov S.K. {\it $\cal P$-differentiability on Carnot groups


in different topologies and related topics} // Тр. по анализу и


геометрии --


Новосибирск. Изд-во Ин-та математики им. С.Л.\,Соболева СО


РАН. -- 2000. -- С.\,603--670.





\bibitem{30}


Водопьянов С.К. {\it Операторы подстановки пространств


Соболева} // Современ. пробл. теории функций и их


прилож., Саратов, янв. 2002 г. -- Саратов. -- 2002. --


С.\,42--43.





\bibitem{31}


Гусман М. {\it Дифференцирование интегралов в ${\Bbb R}^n$.} --


М.: Мир, 1978. -- 200~с.





\bibitem{32}


Херрон А.Д., Коскела П. {\it Области Пуанкаре и


квазиконформные отображения} // Сиб. матем. журн. -- 1995.


-- Т.\,36. -- \No\,6. -- С.\,1416--1434.





\bibitem{33}


Ухлов А.Д. {\it Пространства Соболева и дифференциальные


свойства гомеоморфизмов на группах Карно}. -- Препринт.


Хабаровск. ИПМ ДВО РАН. -- 2000. -- \No\,18. -- 28~с.





\bibitem{34}


Федерер Г. {\it Геометрическая теория меры.} -- М.: Наука,


1987. -- 760~с.





\end{thebibliography}





\vskip 2cm





\post{\it Новосибирский государственный}{\it Поступила}


\post{\it университет}{26.06.2002}





\label{end}


